\chapter{Колокализация компонент спектрального класса пятнадцать}
\label{ch:v6-spectral-transition-component-colocalization-gate}

\section{Критерий связанной частицы}

Классы двенадцать и три существуют как самостоятельные подциклы, а их
сумма восстанавливает глобальный класс пятнадцать. Но топологическое
сложение не является физическим связыванием. Чтобы назвать сумму одной
локализованной частицей, текущий родитель должен содержать хотя бы один из
следующих механизмов:
\begin{enumerate}
 \item прямой оператор, соединяющий кварковую и лептонную опоры;
 \item перекрёстную энергию, зависящую от взаимного положения их ядер;
 \item единственное общее поле массы, которое родительски вынуждает оба
 индекса локализоваться на одном дефекте.
\end{enumerate}

Общий фон или одинаковая нормировка недостаточны. Требуется минимум
энергии при совпадении центров, а не только возможность вручную положить
два профиля в одну точку.

\section{Конечный граф не содержит кварк-лептонного ребра}

Физический модуль одноформ содержит ровно три ребра
\begin{equation}
 u:\ Q_L\leftrightarrow u_R,
 \qquad
 d:\ Q_L\leftrightarrow d_R,
 \qquad
 e:\ L_L\leftrightarrow e_R.
\end{equation}
Их алгебра эндоморфизмов равна $\mathbb C^3$, а внедиагональные
морфизмы между неэквивалентными рёбрами нулевые. KO6 соединяет каждое
ребро с сопряжённым, но не отождествляет $u,d,e$.

Следовательно, конечный оператор Дирака не содержит прямого перехода
между $H_q$ и $H_\ell$. В диаграммном языке конечных спектральных троек
такое отсутствие ребра является настоящим структурным отсутствием
юкавского коннектора, а не выбором базиса
\cite{colocalization-krajewski1998,
colocalization-chamseddine-connes-marcolli2007}.

\section{Моритова кривизна остаётся прямой суммой}

Прямоугольный носитель раскладывается по столбцам:
\begin{equation}
 M_{20\times15}(\mathbb C)
 =M_{20\times12}(\mathbb C)\oplus M_{20\times3}(\mathbb C).
\end{equation}
Для ортогональных столбцовых опор норма относительной кривизны
аддитивна:
\begin{equation}
 \|\mathcal R_q\oplus\mathcal R_\ell\|_{\mathrm{HS}}^2
 =\|\mathcal R_q\|_{\mathrm{HS}}^2
 +\|\mathcal R_\ell\|_{\mathrm{HS}}^2.
\end{equation}
Численный остаток равен $5.69\cdot10^{-14}$ на ненормированном случайном
тесте. Центрированная следовая формула Тома V также содержит сумму двух
угловых норм, а не кварк-лептонный член. Общее левое действие $M_{20}$
является общим чтением, но не оператором расстояния между двумя
пространственными ядрами.

\section{Общий Хиггс не создаёт требуемого коннектора}

Рёбра $d$ и $e$ используют один хиггсовский дублет, но след конечного
оператора остаётся блочным. Спектральное кручение содержит комбинацию
\begin{equation}
 I_{de}=3b^2+c^2,
 \qquad
 b=|y_d|^2,
 \qquad
 c=|y_e|^2,
\end{equation}
и её смешанная производная по $b,c$ равна нулю.

Перекрёстный член $6b^2c^2$ возникает только после дополнительного
возведения $I_{de}$ в квадрат внутри потенциала кручения. Предыдущий гейт
уже запретил этот шаг: такой потенциал является новым функционалом и не
выводится из родителя.

Отдельно минимальный $M_{300}$-аудит дал нулевой коэффициент портала
\begin{equation}
 \Tr(Q^2)H^\dagger H.
\end{equation}
Существующее радиальное сцепление с $|H|^2$ относится к общей амплитуде
семейного состояния и не задаёт относительное положение спектральных
опор $q_q,q_\ell$. Поэтому общий Хиггс не является выведенным механизмом
колокализации.

\section{Семейная связность и поля порядка}

Единая семейная связность действует на фермионах как
\begin{equation}
 A_{\mathrm{fam}}\otimes I_{H_{15}}.
\end{equation}
Она точно коммутирует с $P_q$ и $P_\ell$, поэтому видит обе компоненты
одинаково и не создаёт относительной координаты. Более того, её
динамическая локализация текущим родителем не выведена.

Поле порядка
\begin{equation}
 Q(x)=R(x)-\frac13I_3
\end{equation}
является нейтральным семейным квинтетом $(1,1,0)$. Один профиль $Q(x)$
мог бы условно быть общим фоном для всех пятнадцати состояний. Но его
фермионное сцепление требует не менее пяти независимых видовых
коэффициентов; обычная внутренняя флуктуация не порождает это сцепление,
а минимальный родитель не выбирает единственный ненулевой коэффициент.
Поэтому стабильный бозонный дефект $Q,T,B$ нельзя автоматически
отождествить с общим ядром наблюдаемых компонент.

\section{Единственный условный общий центр}

Каллиасов кандидат имеет массу
\begin{equation}
 \Phi(x)=n(x)\cdot\sigma\otimes(q_q+q_\ell).
 \label{eq:v6-colocalization-common-callias}
\end{equation}
Если такой оператор задан одним профилем $n(x)$, все двенадцать и три
коэффициентные моды действительно локализуются на одном солитонном ядре.
Это стандартная логика фермионной локализации общей знакопеременной или
асимптотически обратимой массой
\cite{colocalization-jackiw-rebbi1976,colocalization-callias1978}.

Но проектный ledger запрещает повысить это условное свойство до вывода:
\begin{enumerate}
 \item пространственный оператор Каллиаса не отождествлён с тёплицевым
 оператором на аналитическом уровне;
 \item конечный родитель не выводит необходимый независимый
 spin-cover-дублет;
 \item двухкопийная лазейка закрыта;
 \item компонентные циклы допускают независимые профили
 $n_q(x),n_\ell(x)$ без нарушения известных симметрий.
\end{enumerate}
Следовательно, формула \eqref{eq:v6-colocalization-common-callias} является
единственным найденным общим ядром, но остаётся недоказанным анзацем.

\section{Прямой тест разделения центров}

Для явного стоп-теста введены две нормированные плотности одинаковой
ширины $\rho_q(x-X_q)$ и $\rho_\ell(x-X_\ell)$. Текущее аддитивное
действие имеет вид
\begin{equation}
 E_0(d)=\frac4{35}\int\rho_q^2dx
       +\frac1{35}\int\rho_\ell^2dx,
 \qquad
 d=|X_q-X_\ell|.
\end{equation}
При $d=0,0.5,1,2,4$ получено одно значение
\begin{equation}
 E_0=0.05037406996
\end{equation}
с разбросом $3.47\cdot10^{-17}$. Энергия не видит расстояние между
центрами.

Для сравнения, гипотетический привлекательный член
\begin{equation}
 -\lambda\int\rho_q\rho_\ell dx,
 \qquad \lambda>0,
\end{equation}
имел бы строгий минимум при $d=0$. В текущем родителе
$\lambda=0$. Это показывает точную недостающую структуру, не вводя её.

\begin{theorem}[Отсутствие выведенной колокализации]
Ни конечный оператор Дирака, ни моритова кривизна, ни допустимое
спектральное кручение, ни минимальный хиггсовский портал, ни семейная
связность, ни поля порядка $Q,T,B$ не создают выведенной перекрёстной
энергии между спектральными классами двенадцать и три. Единый
Каллиасов профиль условно колокализовал бы их, но его физический носитель
текущим родителем не выведен. Поэтому класс пятнадцать не является
доказанной связанной частицей.
\end{theorem}

\begin{remark}
Сохранение полного KO6-класса $15$ допускает пространственно разделённую
сумму зарядов $12+3$. Закон сохранения индекса не является потенциальной
энергией и не требует совпадения центров.
\end{remark}

\section{Следующий различающий тест}

Поскольку связывание полного поколения не найдено, следующий гейт о
хиггсовски разрешённых опорах должен спуститься внутрь компонент и
проверить физические опоры после
выбора ненулевого направления Хиггса. Предварительная, но ещё не
утверждённая ведомость имеет вид
\begin{equation}
 15\stackrel{H\neq0}{\longrightarrow}6_u+6_d+2_e+1_\nu.
\end{equation}
Требуется проверить gauge-ковариантные проекторы, Real-удвоение,
тёплицевы классы и отличие массивных дираковских пар от нейтринной
односторонней линии.

\begin{protocol}
\begin{enumerate}
 \item прямое кварк-лептонное ребро: \textbf{отсутствует};
 \item смешанная моритова норма: \textbf{ноль};
 \item допустимый хиггсовский или торсионный связывающий член:
 \textbf{не выведен};
 \item семейная связность различает опоры $12$ и $3$: \textbf{нет};
 \item поле $Q$ имеет выведенное наблюдаемое фермионное сцепление:
 \textbf{нет};
 \item единый Каллиасов профиль условно даёт общий центр: \textbf{да};
 \item spin-cover-носитель этого профиля выведен: \textbf{нет};
 \item энергия текущего действия зависит от разделения центров:
 \textbf{нет};
 \item глобальный класс $15$ является связанной частицей: \textbf{нет};
 \item следующий тест: \textbf{хиггс-разрешённые минимальные опоры}.
\end{enumerate}
\end{protocol}

Машинный сертификат сохранён в
\path{s2t_v6_spectral_transition_component_colocalization_gate_results.json}.

\begin{thebibliography}{9}
\bibitem{colocalization-krajewski1998}
T.~Krajewski,
\emph{Classification of Finite Spectral Triples},
Journal of Geometry and Physics 28 (1998), 1--30;
arXiv:hep-th/9701081.

\bibitem{colocalization-chamseddine-connes-marcolli2007}
A.~H.~Chamseddine, A.~Connes, M.~Marcolli,
\emph{Gravity and the Standard Model with Neutrino Mixing},
Advances in Theoretical and Mathematical Physics 11 (2007), 991--1089;
arXiv:hep-th/0610241.

\bibitem{colocalization-jackiw-rebbi1976}
R.~Jackiw, C.~Rebbi,
\emph{Solitons with Fermion Number $1/2$},
Physical Review D 13 (1976), 3398--3409.

\bibitem{colocalization-callias1978}
C.~Callias,
\emph{Axial Anomalies and Index Theorems on Open Spaces},
Communications in Mathematical Physics 62 (1978), 213--234.
\end{thebibliography}